\chapter{Trigonometri: Lingkaran Satuan}\label{ch:g11:trigo}

Trigonometri bermula pada segitiga siku-siku, tetapi rumah sejatinya adalah
\emph{lingkaran satuan}, tempat \omterm{def:g11:trigo:cossin}{kosinus} dan \omterm{def:g11:trigo:cossin}{sinus} menjadi \omterm{def:g11:func:function}{fungsi} dari sebarang
\omterm{def:g10:numbers:sets}{bilangan real}. Bab ini memasang \omterm{def:g11:trigo:radian}{radian}, pelilitan garis
\omterm{def:g10:numbers:sets}{bilangan real} pada lingkaran, nilai eksak yang harus dihafal, dan
kesimetrian yang membangkitkan semua identitas klasiknya. Telaah
$\cos$ dan $\sin$ sebagai \omterm{def:g11:func:function}{fungsi} --- variasi, \omterm{def:g11:deriv:derivative}{turunan} --- dilakukan
pada \cref{ch:g12:trigo}.

\section{Radian dan pelilitan}

\begin{definition}[Radian]\label{def:g11:trigo:radian}
Misalkan $\mathcal C$ lingkaran berjari-jari $1$ yang berpusat di titik asal
$O$, dan $I(1, 0)$. Ukuran \emph{radian}\index{radian} sebuah sudut pusat
$\mathcal C$ adalah panjang busur yang dipotongnya pada $\mathcal C$.
Karena keliling lingkaran penuhnya $2\pi$:
\[
360^\circ = 2\pi \text{ rad},
\qquad
180^\circ = \pi \text{ rad},
\qquad
1 \text{ rad} = \frac{180^\circ}{\pi} \approx 57.3^\circ .
\]
\end{definition}

\begin{method}[Mengonversi]\label{met:g11:trigo:convert}
Derajat dan \omterm{def:g11:trigo:radian}{radian} sebanding: kalikan dengan $\frac{\pi}{180}$ untuk beralih
dari derajat ke \omterm{def:g11:trigo:radian}{radian}, dan dengan $\frac{180}{\pi}$ untuk arah sebaliknya.
Misalnya $60^\circ = 60 \times \frac{\pi}{180} = \frac{\pi}{3}$, dan
$\frac{3\pi}{4} = \frac{3 \times 180^\circ}{4} = 135^\circ$.
\end{method}

\begin{definition}[Melilitkan garis pada lingkaran]\label{def:g11:trigo:winding}
Kepada setiap \omterm{def:g10:numbers:sets}{bilangan real} $t$, kaitkan titik $M(t)$ pada $\mathcal C$
yang diperoleh dengan berjalan sejauh $\abs t$ sepanjang lingkaran dari $I$,
berlawanan arah jarum jam bila $t \geq 0$, dan searah jarum jam bila $t < 0$. Karena
kelilingnya $2\pi$, bilangan $t$ dan $t + 2k\pi$ ($k \in \Z$)
mencapai titik yang sama: $M(t + 2k\pi) = M(t)$.
\end{definition}

\begin{example}
$M(0) = I$; $M\!\left(\frac{\pi}{2}\right) = (0, 1)$, yaitu puncak
lingkarannya; $M(\pi) = (-1, 0)$; $M\!\left(-\frac{\pi}{2}\right) = (0, -1)$;
lalu $M\!\left(\frac{9\pi}{4}\right) = M\!\left(\frac{\pi}{4} +
2\pi\right) = M\!\left(\frac{\pi}{4}\right)$.
\end{example}

\section{Kosinus dan sinus}

\begin{definition}[Kosinus dan sinus]\label{def:g11:trigo:cossin}
Untuk setiap \omterm{def:g10:numbers:sets}{bilangan real} $t$, \emph{kosinus}\index{kosinus} dan
\emph{sinus}\index{sinus} dari $t$ adalah \omterm{def:g10:coordgeom:system}{koordinat} titik $M(t)$ pada
lingkaran satuan:
\[
M(t) = (\cos t,\ \sin t).
\]
\end{definition}

\begin{omfigure}
\begin{tikzpicture}[scale=2.0]
  \draw[->] (-1.3,0) -- (1.4,0) node[right] {$x$};
  \draw[->] (0,-1.3) -- (0,1.4) node[above] {$y$};
  \draw[gray] (0,0) circle (1);
  \coordinate (O) at (0,0);
  \coordinate (M) at (55:1);
  \draw[thick,omDef] (O) -- (M);
  \draw[dashed] (M) -- (M |- O) node[below, font=\small] {$\cos t$};
  \draw[dashed] (M) -- (M -| O) node[left, font=\small] {$\sin t$};
  \draw[->,omThm,thick] (0.35,0) arc (0:55:0.35);
  \node[omThm, font=\small] at (27:0.55) {$t$};
  \fill (M) circle (0.7pt) node[above right] {$M(t)$};
  \fill (1,0) circle (0.7pt) node[below right] {$I$};
\end{tikzpicture}

{\small Lingkaran satuan: \omterm{def:g10:numbers:sets}{bilangan real} $t$, yang dililitkan berlawanan arah
jarum jam dari $I$, mendarat di titik $M(t)$ yang \omterm{def:g10:coordgeom:system}{koordinatnya}
\emph{mendefinisikan} $\cos t$ dan $\sin t$.}
\end{omfigure}

\begin{proposition}[Sifat pertama]\label{prop:g11:trigo:first}
Untuk semua $t \in \R$ dan $k \in \Z$:
\[
-1 \leq \cos t \leq 1, \qquad
-1 \leq \sin t \leq 1, \qquad
\cos^2 t + \sin^2 t = 1,
\]
\[
\cos(t + 2k\pi) = \cos t, \qquad \sin(t + 2k\pi) = \sin t .
\]
\end{proposition}

\begin{proof}
$M(t)$ terletak pada lingkaran satuan, sehingga \omterm{def:g10:coordgeom:system}{koordinatnya} berada antara
$-1$ dan $1$, dan \omterm{def:g10:coordgeom:system}{koordinat} itu memenuhi \omterm{def:g10:algebra:equation}{persamaan} lingkarannya
$x^2 + y^2 = 1$, yaitu identitas $\cos^2 t + \sin^2 t = 1$. Kaperiodikannya
menyatakan ulang $M(t + 2k\pi) = M(t)$ (\cref{def:g11:trigo:winding}).
\end{proof}

\begin{example}\label{ex:g11:trigo:identity}
Jika $\sin t = \frac35$ dan $t \in \intcc{\frac\pi2}{\pi}$ (kuadran
kedua), maka $\cos^2 t = 1 - \frac{9}{25} = \frac{16}{25}$, sehingga
$\cos t = \pm\frac45$; pada kuadran kedua absisnya negatif,
jadi $\cos t = -\frac45$.
\end{example}

Nilai yang harus dihafal:
\begin{center}
\begin{tabular}{c|ccccc}
$t$ & $0$ & $\dfrac{\pi}{6}$ & $\dfrac{\pi}{4}$ & $\dfrac{\pi}{3}$ & $\dfrac{\pi}{2}$\\[6pt]
\hline\rule{0pt}{14pt}%
$\cos t$ & $1$ & $\dfrac{\sqrt3}{2}$ & $\dfrac{\sqrt2}{2}$ & $\dfrac12$ & $0$\\[6pt]
$\sin t$ & $0$ & $\dfrac12$ & $\dfrac{\sqrt2}{2}$ & $\dfrac{\sqrt3}{2}$ & $1$
\end{tabular}
\end{center}

\begin{remark}
Kiat mengingatnya: sinusnya terbaca
$\frac{\sqrt0}{2}, \frac{\sqrt1}{2}, \frac{\sqrt2}{2}, \frac{\sqrt3}{2},
\frac{\sqrt4}{2}$, dan kosinusnya adalah daftar yang sama dalam urutan terbalik.
\end{remark}

\section{Sudut berelasi}

\begin{proposition}[Sudut berelasi]\label{prop:g11:trigo:associated}
\index{sudut berelasi}
Untuk semua $t \in \R$:
\[
\begin{aligned}
\cos(-t) &= \cos t, & \sin(-t) &= -\sin t,\\
\cos(\pi - t) &= -\cos t, & \sin(\pi - t) &= \sin t,\\
\cos(\pi + t) &= -\cos t, & \sin(\pi + t) &= -\sin t,\\
\cos\left(\tfrac{\pi}{2} - t\right) &= \sin t, &
\sin\left(\tfrac{\pi}{2} - t\right) &= \cos t .
\end{aligned}
\]
\end{proposition}

\begin{proof}
Setiap barisnya menyatakan sebuah kesimetrian lingkaran.

$M(-t)$ adalah cerminan $M(t)$ terhadap sumbu $x$: absisnya
tidak berubah, \omterm{def:g10:coordgeom:system}{ordinatnya} berganti tanda.

$M(\pi - t)$ adalah cerminan $M(t)$ terhadap sumbu $y$: berjalan
mundur dari $\pi$ sejauh $t$ mendarat berseberangan (secara mendatar) dengan
berjalan maju dari $0$ sejauh $t$. Absisnya berganti tanda, \omterm{def:g10:coordgeom:system}{ordinatnya}
tidak berubah.

$M(\pi + t)$ terletak diametral berseberangan dengan $M(t)$ (setengah putaran):
kedua \omterm{def:g10:coordgeom:system}{koordinatnya} berganti tanda.

$M\!\left(\frac\pi2 - t\right)$ adalah cerminan $M(t)$ terhadap
garis diagonal $y = x$, yang menukar kedua \omterm{def:g10:coordgeom:system}{koordinatnya}.
\end{proof}

\begin{omfigure}
\begin{tikzpicture}[scale=2.0]
  \draw[->] (-1.35,0) -- (1.45,0) node[right] {$x$};
  \draw[->] (0,-1.35) -- (0,1.45) node[above] {$y$};
  \draw[gray] (0,0) circle (1);
  \fill (40:1) circle (0.7pt) node[above right, font=\small] {$M(t)$};
  \fill (-40:1) circle (0.7pt) node[below right, font=\small] {$M(-t)$};
  \fill (140:1) circle (0.7pt) node[above left, font=\small]
    {$M(\pi - t)$};
  \fill (220:1) circle (0.7pt) node[below left, font=\small]
    {$M(\pi + t)$};
  \draw[black, dashed, thin] (40:1) -- (-40:1);
  \draw[black, dashed, thin] (40:1) -- (140:1);
  \draw[black, dashed, thin] (140:1) -- (220:1);
  \draw[omDef, thin] (0,0) -- (40:1);
  \draw[omDef, thin] (0,0) -- (140:1);
  \draw[omDef, thin] (0,0) -- (220:1);
  \draw[omDef, thin] (0,0) -- (-40:1);
\end{tikzpicture}

{\small Keempat titik berelasinya: $M(-t)$ mencerminkan $M(t)$ terhadap
sumbu $x$, $M(\pi - t)$ terhadap sumbu $y$, dan $M(\pi + t)$ terletak diametral
berseberangan. Setiap identitas pada \cref{prop:g11:trigo:associated} terbaca
dari gambar ini.}
\end{omfigure}

\begin{example}
$\cos\frac{2\pi}{3} = \cos\left(\pi - \frac\pi3\right)
= -\cos\frac\pi3 = -\frac12$, dan
$\sin\frac{7\pi}{6} = \sin\left(\pi + \frac\pi6\right)
= -\sin\frac\pi6 = -\frac12$. Tabel berisi lima nilai itu, yang diperluas oleh
kesimetriannya, meliputi seluruh lingkaran.
\end{example}

\section{Menyelesaikan persamaan trigonometri}

\begin{theorem}[Persamaan $\cos t = \cos a$ dan $\sin t = \sin a$]
\label{thm:g11:trigo:equations}
Untuk \omterm{def:g10:numbers:sets}{bilangan real} $t$ dan $a$:
\[
\cos t = \cos a \iff t = a + 2k\pi \ \text{atau}\ t = -a + 2k\pi
\quad (k \in \Z),
\]
\[
\sin t = \sin a \iff t = a + 2k\pi \ \text{atau}\ t = \pi - a + 2k\pi
\quad (k \in \Z).
\]
\end{theorem}

\begin{proof}
Dua titik pada lingkaran satuan berabsis sama tepat ketika keduanya
berimpit atau saling mencerminkan terhadap sumbu $x$; menurut
\cref{prop:g11:trigo:associated}, cerminan $M(a)$ adalah $M(-a)$. Jadi
$\cos t = \cos a$ berarti $M(t) = M(a)$ atau $M(t) = M(-a)$, yaitu
$t = \pm a$ sampai pada kelipatan $2\pi$. Serupa dengan itu, \omterm{def:g10:coordgeom:system}{ordinat} yang sama
berarti titiknya berimpit atau saling mencerminkan terhadap sumbu $y$, dan
cerminan $M(a)$ adalah $M(\pi - a)$.
\end{proof}

\begin{method}[Menyelesaikan $\cos t = c$ pada sebuah interval]\label{met:g11:trigo:solve}
\begin{enumerate}
  \item Carilah \emph{satu} sudut $a$ dengan $\cos a = c$, dari tabel
        nilai yang sudah dikenal.
  \item Tulislah penyelesaian umumnya $t = \pm a + 2k\pi$
        (\cref{thm:g11:trigo:equations}).
  \item Pilihlah \omterm{def:g10:numbers:sets}{bilangan bulat} $k$ yang mendarat di dalam \omterm{def:g10:numbers:interval}{interval} yang
        diminta, lalu daftarkan penyelesaiannya.
\end{enumerate}
Lakukan hal serupa untuk $\sin t = c$ dengan $t = a + 2k\pi$ atau
$t = \pi - a + 2k\pi$.
\end{method}

\begin{example}\label{ex:g11:trigo:solve}
Selesaikan $\cos t = \frac12$ pada $\intoc{-\pi}{\pi}$. Satu penyelesaian
$\cos a = \frac12$ adalah $a = \frac{\pi}{3}$. Penyelesaian umumnya adalah
$t = \frac\pi3 + 2k\pi$ dan $t = -\frac\pi3 + 2k\pi$. Di dalam
$\intoc{-\pi}{\pi}$, hanya $k = 0$ yang menyumbang:
$t \in \left\{-\frac\pi3,\ \frac\pi3\right\}$.

Selesaikan $\sin t = -\frac{\sqrt2}{2}$ pada $\intco{0}{2\pi}$. Satu sudutnya
adalah $a = -\frac\pi4$; penyelesaiannya $t = -\frac\pi4 + 2k\pi$ dan
$t = \pi + \frac\pi4 + 2k\pi = \frac{5\pi}{4} + 2k\pi$. Di
$\intco{0}{2\pi}$: $t = \frac{7\pi}{4}$ (dari $k = 1$) dan
$t = \frac{5\pi}{4}$.
\end{example}

\begin{omfigure}
\begin{tikzpicture}[scale=2.0]
  \draw[->] (-1.35,0) -- (1.45,0) node[right] {$x$};
  \draw[->] (0,-1.2) -- (0,1.3) node[above] {$y$};
  \draw[gray] (0,0) circle (1);
  \draw[omProp, thick] (0.5,-1.05) -- (0.5,1.15)
    node[above, font=\small] {$x = \tfrac12$};
  \fill (60:1) circle (0.7pt)
    node[above right, font=\small] {$M\!\left(\tfrac\pi3\right)$};
  \fill (-60:1) circle (0.7pt)
    node[below right, font=\small] {$M\!\left(-\tfrac\pi3\right)$};
  \draw[omDef, thin] (0,0) -- (60:1);
  \draw[omDef, thin] (0,0) -- (-60:1);
\end{tikzpicture}

{\small Menyelesaikan $\cos t = \frac12$: garis tegak $x = \frac12$
(jingga) memotong lingkaran satuan di dua titik, yang setangkup terhadap
sumbu $x$ --- karena itulah kedua keluarga penyelesaiannya $t = \pm\frac\pi3 + 2k\pi$.}
\end{omfigure}

\section{Latihan}

\begin{exercise}[$\star$]\label{exo:g11:trigo:1}
Konversikan ke \omterm{def:g11:trigo:radian}{radian}: $30^\circ$, $45^\circ$, $120^\circ$, $270^\circ$.
Konversikan ke derajat: $\frac{\pi}{5}$, $\frac{5\pi}{6}$, $\frac{7\pi}{4}$,
$\frac{2\pi}{9}$.
\end{exercise}

\begin{exercise}[$\star$]\label{exo:g11:trigo:2}
Tempatkan pada lingkaran satuan titik $M(t)$ untuk
\[
t = \frac{2\pi}{3},\quad
t = -\frac{\pi}{4},\quad
t = \frac{17\pi}{6},\quad
t = -\frac{7\pi}{2},
\]
setelah masing-masing disusutkan ke sebuah nilai di $\intoc{-\pi}{\pi}$ modulo $2\pi$.
\end{exercise}

\begin{exercise}[$\star$]\label{exo:g11:trigo:3}
Dengan memakai tabel dan \omterm{prop:g11:trigo:associated}{sudut berelasi}, berikan nilai eksak
\[
\cos\frac{3\pi}{4}, \quad
\sin\frac{5\pi}{6}, \quad
\cos\left(-\frac{\pi}{3}\right), \quad
\sin\frac{4\pi}{3}, \quad
\cos\frac{11\pi}{6} .
\]
\end{exercise}

\begin{exercise}[$\star$]\label{exo:g11:trigo:4}
Diberikan $\cos t = \frac{5}{13}$ dengan $t \in \intoo{-\frac\pi2}{0}$,
hitunglah $\sin t$ secara eksak.
\end{exercise}

\begin{exercise}[$\star\star$]\label{exo:g11:trigo:5}
Sederhanakan, untuk sebarang \omterm{def:g10:numbers:sets}{bilangan real} $x$:
\[
A = \cos(\pi + x) + \cos(-x) + \sin\left(\frac\pi2 - x\right) - \cos x,
\qquad
B = \sin(\pi - x) + \sin(\pi + x) + \sin(-x) .
\]
\end{exercise}

\begin{exercise}[$\star\star$]\label{exo:g11:trigo:6}
Selesaikan pada $\intoc{-\pi}{\pi}$:
\[
\cos t = \frac{\sqrt3}{2}, \qquad
\sin t = \frac12, \qquad
\cos t = -1 .
\]
\end{exercise}

\begin{exercise}[$\star\star$]\label{exo:g11:trigo:7}
Selesaikan pada $\intco{0}{2\pi}$:
\[
\sin t = -\frac{\sqrt3}{2}, \qquad
\cos t = 0, \qquad
2\sin t + 1 = 0 .
\]
\end{exercise}

\begin{exercise}[$\star\star$]\label{exo:g11:trigo:8}
Selesaikan pada $\R$ \omterm{def:g10:algebra:equation}{persamaan} $\cos 2t = \cos t$. (Pakailah
\cref{thm:g11:trigo:equations} dengan $a = t$, lalu bahaslah kedua keluarga
penyelesaiannya.)
\end{exercise}

\begin{exercise}[$\star\star$]\label{exo:g11:trigo:9}
Tunjukkan bahwa untuk semua \omterm{def:g10:numbers:sets}{bilangan real} $x$ berlaku
\[
\bigl(\cos x + \sin x\bigr)^2 + \bigl(\cos x - \sin x\bigr)^2 = 2 .
\]
\end{exercise}

\begin{exercise}[$\star\star$]\label{exo:g11:trigo:10}
Sebuah roda berjari-jari $30$ cm menggelinding tanpa tergelincir. Sudut berapa
(dalam \omterm{def:g11:trigo:radian}{radian}, lalu dalam derajat) yang ditempuhnya ketika sepedanya maju
$1.5$ m? Sejauh apa sepeda itu maju selama satu putaran penuh rodanya?
\end{exercise}

\begin{exercise}[$\star\star\star$]\label{exo:g11:trigo:11}
Selesaikan pada $\intoc{-\pi}{\pi}$ pertidaksamaan
$\cos t \leq \frac12$. (Sketsalah lingkaran satuannya, tandai daerah yang
absisnya paling besar $\frac12$, lalu bacalah \omterm{def:g10:numbers:interval}{interval} nilai
$t$-nya.)
\end{exercise}

\section{Soal: Kincir ria dan pasang surut}

\begin{problem}[{Soal akhir pekan --- radian sebagai jarak gelinding,
dan lingkaran satuan sebagai jam induk setiap gejala
berkala}]
\label{pb:g11:trigo:1}
Sebuah kincir ria membawamu mengelilingi lingkaran; pasang surut membawa
air pelabuhan naik dan turun pada lingkaran matematis yang sama. Setiap
gejala yang berkala --- roda, pasang surut, bunyi, musim --- membaca
jadwalnya dari lingkaran satuan bab ini
(\cref{def:g11:trigo:cossin}), dan \omterm{def:g10:algebra:equation}{persamaan} pada
\cref{met:g11:trigo:solve} memberitahukan waktunya secara tepat. Soal ini
mengonversi, melilit, memodelkan dan menyelesaikan --- lalu di sepanjang
jalan itu menuntaskan mengapa matematikawan mengukur sudut dalam \omterm{def:g11:trigo:radian}{radian}.

\medskip
\noindent\textbf{Bagian I --- \omterm{def:g11:trigo:radian}{Radian}, jarak yang tergelinding.}
\begin{enumerate}
  \item Konversikan ke \omterm{def:g11:trigo:radian}{radian}: $30^\circ$, $135^\circ$,
        $300^\circ$; lalu ke derajat: $\frac{3\pi}{4}$,
        $\frac{5\pi}{6}$ (\cref{met:g11:trigo:convert}).
  \item Inti \omterm{def:g11:trigo:radian}{radian}: sebuah busur bersudut $t$ (dalam
        \omterm{def:g11:trigo:radian}{radian}) pada lingkaran berjari-jari $r$ berpanjang tepat
        $r\,t$. Busur berapa yang dipotong sudut $2$ \omterm{def:g11:trigo:radian}{radian} pada
        lingkaran berjari-jari $3$ m? (Tanpa $\pi$ di mana pun --- itulah
        intinya.)
  \item Sebuah roda berjari-jari $30$ cm menggelinding $1$ km tanpa
        tergelincir. Berapa \emph{\omterm{def:g11:trigo:radian}{radian}} yang telah ditempuhnya,
        dan berapa putaran penuh?
  \item Tempatkan pada lingkaran satuan lalu hitunglah secara eksak:
        $\cos\frac{2\pi}{3}$, $\sin\left(-\frac{\pi}{4}\right)$,
        dan $\cos\frac{19\pi}{6}$ (susutkan modulo $2\pi$ lebih dahulu;
        \cref{prop:g11:trigo:associated}).
  \item Sebuah sudut $t \in \intoo{\frac\pi2}{\pi}$ memenuhi
        $\sin t = \frac35$. Dengan memakai
        $\cos^2 t + \sin^2 t = 1$ dan kuadrannya, carilah
        $\cos t$ dan $\tan t$.
\end{enumerate}

\medskip
\noindent\textbf{Bagian II --- Kincir ria.}
Sebuah kincir raksasa berjari-jari $60$ m, dengan porosnya $65$ m di atas
tanah, dan berputar sekali setiap $30$ menit. Kamu naik di
titik terendah pada waktu $t = 0$ (dalam menit), sehingga tinggimu menuruti
\[
h(t) = 65 - 60 \cos\!\left(\frac{\pi t}{15}\right).
\]
\begin{enumerate}[resume]
  \item Berikan alasan rumusnya: periksalah sudut yang ditempuh setelah $t$
        menit, tingginya di $t = 0$, lalu jelaskan tanda minusnya.
  \item Hitunglah tinggimu pada $t = 7.5$, $t = 15$ dan
        $t = 20$ menit (nilai eksaknya).
  \item Pada waktu berapa kamu \emph{pertama kali} berada di $95$ m?
        (Selesaikan $h(t) = 95$ dengan \cref{met:g11:trigo:solve}.)
  \item Pada selang waktu mana selama perjalanan itu kamu berada paling
        sedikit $95$ m di atas tanah --- dan berapa lama seluruhnya?
        (Bandingkan \cref{exo:g11:trigo:11}.)
  \item Di mana tingginya berubah \emph{paling cepat}? Bandingkan
        kenaikan selama menit pertama
        ($h(1) - h(0)$) dengan kenaikan antara menit $7$ dan
        $8$, lalu rumuskan pengamatannya (perkakas yang membuatnya
        eksak --- \omterm{def:g11:deriv:derivative}{turunan} sinusoida --- tiba
        pada tahun berikutnya).
  \item Giliran sang insinyur: rancanglah rumus tinggi sebuah kincir
        dengan tinggi \omterm{def:g10:functions:extrema}{maksimum} $40$ m, tinggi \omterm{def:g10:functions:extrema}{minimum} $4$ m dan
        kala $12$ menit, dengan penumpang naik di bawah pada
        $t = 0$.
\end{enumerate}

\medskip
\noindent\textbf{Bagian III --- Pasang surut dan kesimetrian
lingkaran.}
Di sebuah pelabuhan, kedalaman airnya dalam meter, $t$ jam setelah
tengah malam, adalah
\[
d(t) = 7 + 3 \sin\!\left(\frac{\pi t}{6}\right).
\]
\begin{enumerate}[resume]
  \item Berikan kedalaman \omterm{def:g10:functions:extrema}{maksimum} dan \omterm{def:g10:functions:extrema}{minimumnya}, waktu keduanya
        pertama kali terjadi, dan kala pasang surutnya.
  \item Sebuah kapal barang memerlukan $8.5$ m air. Selesaikan
        $d(t) \geq 8.5$ pada $\intco{0}{12}$: pada jendela waktu mana
        kapal itu dapat masuk, dan berapa lama jendela itu? Kapan
        jendela berikutnya terbuka?
  \item Turunkan ulang dua identitas \omterm{prop:g11:trigo:associated}{sudut berelasi} langsung dari
        kesimetrian lingkaran --- $\sin(\pi - t) = \sin t$ (cerminan
        terhadap sumbu tegak) dan $\cos(\pi + t) = -\cos t$
        (setengah putaran) --- lalu nyatakan identitas pelengkapnya
        $\cos\left(\frac\pi2 - t\right) = \sin t$.
  \item Paritas (dengan kosakata \cref{pb:g11:func:1}): golongkan
        $\sin$, $\cos$, dan $t \mapsto \sin t + t^3$ sebagai genap atau
        ganjil, dengan bukti satu baris dari lingkarannya.
  \item Selesaikan pada $\intco{0}{2\pi}$:
        $2\sin^2 t - \sin t - 1 = 0$. (Sebuah bentuk kuadrat dalam
        samaran: faktorkan seperti pada \cref{pb:g11:quad:1}, lalu
        bacalah lingkarannya.)
\end{enumerate}

\medskip
\noindent\textbf{Bagian IV --- Sudut milik matematikawan.}
\begin{enumerate}[resume]
  \item Berapa derajat $1$ \omterm{def:g11:trigo:radian}{radian}? Lalu sebuah jebakan kalkulator:
        mana yang lebih besar, $\sin(1)$ (mode \omterm{def:g11:trigo:radian}{radian}) atau
        $\sin(1^\circ)$ --- dan kira-kira berapa kali lipat? Apa moralnya
        bagi penyetelan kalkulator?
  \item Selesaikan $\cos t = \sin t$ pada $\intco{0}{2\pi}$, dengan
        memakai lingkarannya (di mana \omterm{def:g10:coordgeom:system}{absis} dan \omterm{def:g10:coordgeom:system}{ordinatnya} sama?).
  \item Untuk sudut kecil, busur dan sinusnya hampir berimpit:
        bandingkan $\sin(0.1)$ dengan $0.1$ (lima angka desimal), lalu
        berikan galat nisbinya. ``\omterm{def:g10:numbers:approx}{Hampiran} sudut kecil'' ini
        mengemudikan kapal dan bandul; teorema
        di baliknya ($\frac{\sin t}{t} \to 1$) menjadi bintang bab
        limit tahun berikutnya.
  \item Penutup --- lingkaran satuan itu sebagai jam induk, satu
        kalimat untuk setiap butir: \omterm{def:g11:trigo:radian}{radian} mengubah sudut menjadi bilangan yang jujur
        (busur $=$ jari-jari $\times$ sudut); \omterm{def:g11:trigo:cossin}{kosinus} dan \omterm{def:g11:trigo:cossin}{sinus} adalah
        bayangan dan tinggi jarum jamnya; kesimetrian lingkarannya
        adalah identitasnya; menyelesaikan $\cos t = c$ membaca waktu
        dari jam itu; dan setiap gejala berkala --- roda,
        pasang surut, bunyi --- adalah jam ini yang mengenakan busana.

\end{enumerate}
\end{problem}
